% Included after Appendix A of the English report. Requires amsmath and amssymb.
\numberwithin{equation}{section}
\section{Fourier-space equations and the split-step algorithm}\label{sec:B}

\subsection{What changes under a Fourier transform}\label{sec:B.1}
With the Fourier convention of Eq.~\eqref{eq:32}, spatial differentiation gives
\(\mathcal F[\nabla_\perp^2\psi]=-k^2\widetilde\psi\), where
\(k^2=k_x^2+k_y^2\). Transforming the full equations gives
\begin{align}
\partial_t\widetilde\psi={}&-\frac{iKk^2}{\hbar}\widetilde\psi
-\frac{\gamma_C}{2}\widetilde\psi
+\frac{R}{2}\mathcal F[n_R\psi]\notag\\
&-\frac{i}{\hbar}\left(g_C\mathcal F[|\psi|^2\psi]
+g_R\mathcal F[(n_R+\eta P)\psi]\right),\label{eq:B.field}\\
\partial_t\widetilde n_R={}&\widetilde P-\gamma_R\widetilde n_R
-R\mathcal F[n_R|\psi|^2].\label{eq:B.reservoir}
\end{align}
Only the kinetic term is diagonal in wave vector. Products in real space
become convolutions in Fourier space, so the full evolution does not reduce
to independent equations for each wave vector. The implementation therefore
uses Fourier space for kinetic evolution and real space for the local terms.

\subsection{Kinetic half-steps}\label{sec:B.2}
During the kinetic substep, the field obeys
\begin{equation}
\begin{aligned}
\partial_t\widetilde\psi&=-\frac{iKk^2}{\hbar}\widetilde\psi,\\
\widetilde\psi(t+\tau)&=e^{-iKk^2\tau/\hbar}\widetilde\psi(t),
\qquad \partial_t n_R=0.
\end{aligned}\label{eq:B.kinetic}
\end{equation}
The exact discrete kinetic map is therefore
\begin{equation}
\mathcal T_\tau\psi=
\operatorname{IFFT2}\!\left[
e^{-iK(k_x^2+k_y^2)\tau/\hbar}
\operatorname{FFT2}(\psi)\right].\label{eq:B.kinetic-map}
\end{equation}
Let \(\mathcal L_\tau\) denote the local, coupled field--reservoir
evolution defined below. One implemented Strang step is
\begin{equation}
\begin{aligned}
(\psi^{(1)},n_R^{(1)})&=(\mathcal T_{\Delta t/2}\psi^n,n_R^n),\\
(\psi^{(2)},n_R^{(2)})&=\mathcal L_{\Delta t}
(\psi^{(1)},n_R^{(1)}),\\
(\psi^{n+1},n_R^{n+1})&=(\mathcal T_{\Delta t/2}\psi^{(2)},n_R^{(2)}).
\end{aligned}\label{eq:B.strang}
\end{equation}
Thus, each kinetic half-step performs FFT, multiplication by the phase
factor, and inverse FFT. The ordering is a kinetic half-step, a full local
step, and a second kinetic half-step. It is not a static-potential update:
the local coefficients change with the evolving field and reservoir.

\subsection{Local evolution}\label{sec:B.3}
At each real-space node, the local substep solves
\begin{equation}
\begin{aligned}
\dot\psi&=\left[\frac{Rn_R-\gamma_C}{2}
-\frac{i}{\hbar}\bigl(g_C|\psi|^2+g_R(n_R+\eta P)\bigr)\right]\psi,\\
\dot n_R&=P-(\gamma_R+R|\psi|^2)n_R.
\end{aligned}\label{eq:B.local}
\end{equation}
Write this system as \(\dot Y=f(Y;P)\), with
\(Y=(\operatorname{Re}\psi,\operatorname{Im}\psi,n_R)^T\).
For the nonlinear condensate--reservoir calculations, the local map is evaluated by a fourth-order Runge--Kutta step:
\begin{equation}
\begin{aligned}
a_1&=f(Y;P),& a_2&=f(Y+\tfrac{\Delta t}{2}a_1;P),\\
a_3&=f(Y+\tfrac{\Delta t}{2}a_2;P),& a_4&=f(Y+\Delta t\,a_3;P),\\
Y_{\mathrm{out}}&=Y+\frac{\Delta t}{6}(a_1+2a_2+2a_3+a_4).
\end{aligned}\label{eq:B.rk4}
\end{equation}
The density \(n_C=|\psi|^2\) and the reservoir-dependent coefficients
are recomputed at each stage; \(n_C\) is not integrated as a separate
variable. Nodes evolve independently only within this local substep.
The kinetic substeps connect the spatial field. Although the kinetic map
is exact for the discrete Fourier modes and the local solver is fourth
order, the complete splitting scheme is second order in time.


For the linear threshold search, the reservoir is instead frozen at
\(n_{R,0}=P/\gamma_R\), and the cubic field term is omitted, as in
Eq.~\eqref{eq:36}. The local coefficient is then time-independent,
so the local step is evaluated by an exact exponential rather than RK4:
\begin{equation}\label{eq:B.linear-local}
\psi_{\mathrm{out}}=\exp\!\left\{\Delta t\left[
\frac{RP/\gamma_R-\gamma_C}{2}
-\frac{i g_R}{\hbar}\left(\frac{P}{\gamma_R}+\eta P\right)
\right]\right\}\psi_{\mathrm{in}}.
\end{equation}
The same kinetic half-steps surround this local step. The exponential
is exact only for the frozen-reservoir local problem; splitting the
spatial operators still gives a second-order scheme.

\subsection{Spatial grids and their checks}\label{sec:B.4}
The field \(\psi\) and reservoir density \(n_R\) occupy the same uniform,
periodic Cartesian grid. The pump is a fixed array on this grid, and
\(n_C\) is computed pointwise. For the even grid sizes used here, coordinates and FFT wave vectors are
\begin{equation}
\begin{aligned}
x_j&=-\frac{L_x}{2}+j\Delta x,
&\Delta x&=\frac{L_x}{N_x},\quad j=0,\ldots,N_x-1,\\
y_l&=-\frac{L_y}{2}+l\Delta y,
&\Delta y&=\frac{L_y}{N_y},\quad l=0,\ldots,N_y-1,\\
k_{x,m}&=\frac{2\pi m}{L_x},
&m&=0,\ldots,\frac{N_x}{2}-1,-\frac{N_x}{2},\ldots,-1,
\end{aligned}\label{eq:B.grid}
\end{equation}
with the analogous ordering for \(k_y\). The upper spatial endpoint is
not repeated. The FFT normalization cancels between the forward and
inverse transforms in Eq.~\eqref{eq:B.kinetic-map}.

\begin{center}
\small\begin{tabular}{lrrrr}
\hline
Calculation & Grid & \(L\) (\(\mu\mathrm m\)) & \(\Delta x\) (\(\mu\mathrm m\)) & \(\Delta t\) (ps)\\
\hline
Nonlinear spot & \(512^2\) & 256 & 0.5 & 0.025\\
Nonlinear ring & \(512^2\) & 256 & 0.5 & 0.05\\
Linear threshold scan & \(128^2\) & 128 & 1 & 0.05\\
\hline
\end{tabular}
\end{center}
For the main nonlinear runs, the pump scale \(d=2.5\,\mu\mathrm m\)
spans five spatial steps, and the domain is much larger than the pumped
region. These scale comparisons motivate the grid but do not establish
accuracy by themselves. Grid, time-step, boundary, and spectral checks
are summarized in Appendix~\ref{sec:A.4}. The linear threshold endpoints are also
checked on a \(256^2\) grid at \(L=128\,\mu\mathrm m\) and a
\(192^2\) grid at \(L=192\,\mu\mathrm m\). Spatial resolution and
domain size are separate checks.
